% !TEX root = ../notes.tex

\documentclass[../notes.tex]{subfiles}

\begin{document}

\section{April 8}
Today, we continue playing the Vogan game.

\subsection{Real Forms of Types \texorpdfstring{$F_4$}{ F4} and \texorpdfstring{$G_2$}{ G2}}
We now pass to exceptional types. Linear algebra is no longer so helpful to find our real forms. Instead, we will use Vogan diagrams. We have already seen that many Vogan diagrams can give the same real form, so it will be helpful to have a means to determine when two Vogan diagrams give the same real form. Here is one criterion.
\begin{example} \label{ex:vogan-game}
	Choose a Vogan diagram coming from some $\theta$, and choose a vertex $i$ stable under $\theta$. Then $\theta$ and $s_i\theta$ give the same real form because they are automatically conjugate. On Vogan diagrams, one sees that
	\[\alpha_j(s_i\theta)=s_i\alpha_j(\theta)=(\alpha_j-a_{ij}\alpha_i)(\theta)=\alpha_j(\theta)\alpha_i(\theta)^{-a_{ij}}.\]
	Now, $\alpha_j(\theta)$ is $+1$ when $j$ is colored white, and it is $-1$ when $j$ is colored black. Thus, the Vogan diagram is entirely the same, except when $\alpha_i(\theta)=-1$ and $a_{ij}$ is odd. In practice, this means that we can take a black vertex $i$ in the Vogan diagram, and we may change the colors of all neighbors $j$ with $a_{ij}\in\{-1,-3\}$ to receive a Vogan diagram with an equivalent real form.
\end{example}
The above rule allows us to play the ``Vogan game'' to try to go between real forms.\footnote{Because the Weyl group acts transitively on polarizations, this rule should be enough to exhaust the relations in the compact inner class.}
\begin{example}
	In type $G_2$, the Dynkin diagram has no nontrivial automorphisms. The Vogan diagram
	% https://q.uiver.app/#q=WzAsMixbMCwwLCJcXGNpcmMiXSxbMSwwLCJcXGNpcmMiXSxbMSwwLCIiLDAseyJsZXZlbCI6M31dXQ==&macro_url=https%3A%2F%2Fraw.githubusercontent.com%2FdFoiler%2Fnotes%2Fmaster%2Fnir.tex
	% https://q.uiver.app/#q=WzAsMixbMCwwLCJcXGNpcmMiXSxbMSwwLCJcXGNpcmMiXSxbMSwwLCIiLDEseyJsZXZlbCI6Mn1dLFsxLDAsIiIsMSx7InN0eWxlIjp7ImhlYWQiOnsibmFtZSI6Im5vbmUifX19XV0=&macro_url=https%3A%2F%2Fraw.githubusercontent.com%2FdFoiler%2Fnotes%2Fmaster%2Fnir.tex
	\[\begin{tikzcd}[cramped]
		\circ & \circ
		\arrow[Rightarrow, from=1-2, to=1-1]
		\arrow[no head, from=1-2, to=1-1]
	\end{tikzcd}\]
	gives the compact form; here, $\mf k=\mf g$. There is also the (non-isomorphic) split form, which must come from the Vogan diagrams
	% https://q.uiver.app/#q=WzAsNixbMCwwLCJcXGJ1bGxldCJdLFsxLDAsIlxcY2lyYyJdLFsyLDAsIlxcYnVsbGV0Il0sWzMsMCwiXFxidWxsZXQiXSxbNCwwLCJcXGNpcmMiXSxbNSwwLCJcXGJ1bGxldCJdLFszLDIsIiIsMCx7ImxldmVsIjoyfV0sWzUsNCwiIiwwLHsibGV2ZWwiOjJ9XSxbMSwwLCIiLDEseyJsZXZlbCI6Mn1dLFsxLDAsIiIsMSx7InN0eWxlIjp7ImhlYWQiOnsibmFtZSI6Im5vbmUifX19XSxbMywyLCIiLDIseyJzdHlsZSI6eyJoZWFkIjp7Im5hbWUiOiJub25lIn19fV0sWzUsNCwiIiwyLHsic3R5bGUiOnsiaGVhZCI6eyJuYW1lIjoibm9uZSJ9fX1dXQ==&macro_url=https%3A%2F%2Fraw.githubusercontent.com%2FdFoiler%2Fnotes%2Fmaster%2Fnir.tex
	\[\begin{tikzcd}[cramped]
		\bullet & \circ & \bullet & \bullet & \circ & \bullet
		\arrow[Rightarrow, from=1-2, to=1-1]
		\arrow[no head, from=1-2, to=1-1]
		\arrow[Rightarrow, from=1-4, to=1-3]
		\arrow[no head, from=1-4, to=1-3]
		\arrow[Rightarrow, from=1-6, to=1-5]
		\arrow[no head, from=1-6, to=1-5]
	\end{tikzcd}\]
	which are all equivalent (via \Cref{ex:vogan-game}). For the split inner form, one can compute as in \Cref{thm:compact-form-compact} that the signature of the Killing form on the split form is $(\op{rank}\mf g+\#\Phi_+,\#\Phi_+)$, so comparing with the splitting $\mf g_\theta=\mf k^c\oplus i\mf p^c$ shows that the rank of $\mf k$ is $\#\Phi_+$. Thus, $\mf k$ is some six-dimensional Lie algebra of rank $2$; the only possibility is $\mf{sl}(2)\oplus\mf{sl}(2)$.
\end{example}
\begin{example}
	We compute the Vogan diagrams in type $F_4$.
\end{example}
\begin{proof}
	Before starting, recall that $F_4$ is a root system in $\frac12\ZZ^4$, given by the roots in $B_4$ (which are $\{\pm e_i\}\sqcup\{\pm e_i\pm e_j\}$) and adding the roots $\pm_1e_1\pm_2e_2\pm_3e_3\pm_4e_4$. Thus, $F_4$ has rank $4$ and dimension $4+32+16=52$. The Dynkin diagram has no nontrivial automorphisms. The Vogan diagram
	% https://q.uiver.app/#q=WzAsNCxbMCwwLCJcXGNpcmMiXSxbMSwwLCJcXGNpcmMiXSxbMiwwLCJcXGNpcmMiXSxbMywwLCJcXGNpcmMiXSxbMiwxLCIiLDAseyJsZXZlbCI6Mn1dLFszLDIsIiIsMCx7InN0eWxlIjp7ImhlYWQiOnsibmFtZSI6Im5vbmUifX19XSxbMSwwLCIiLDAseyJzdHlsZSI6eyJoZWFkIjp7Im5hbWUiOiJub25lIn19fV1d&macro_url=https%3A%2F%2Fraw.githubusercontent.com%2FdFoiler%2Fnotes%2Fmaster%2Fnir.tex
	\[\begin{tikzcd}[cramped]
		\circ & \circ & \circ & \circ
		\arrow[no head, from=1-2, to=1-1]
		\arrow[Rightarrow, from=1-3, to=1-2]
		\arrow[no head, from=1-4, to=1-3]
	\end{tikzcd}\]
	gives the compact form; here, $\mf k=\mf g$.
	
	We now play the Vogan game. By inspection, we see that $a_{ij}$ is odd for all $i\ne j$, except for $a_{23}=-2$. Thus, the right two vertices (which are the long roots) can change the colors of the left two vertices (which are the short roots), but not the other way around. To this end, note that all the diagrams
	% https://q.uiver.app/#q=WzAsNixbMiwwLCJcXGJ1bGxldCJdLFszLDAsIlxcYnVsbGV0Il0sWzQsMCwiXFxjaXJjIl0sWzUsMCwiXFxidWxsZXQiXSxbMCwwLCJcXGJ1bGxldCJdLFsxLDAsIlxcY2lyYyJdLFswLDEsIiIsMCx7InN0eWxlIjp7ImhlYWQiOnsibmFtZSI6Im5vbmUifX19XSxbNCw1LCIiLDAseyJzdHlsZSI6eyJoZWFkIjp7Im5hbWUiOiJub25lIn19fV0sWzIsMywiIiwwLHsic3R5bGUiOnsiaGVhZCI6eyJuYW1lIjoibm9uZSJ9fX1dXQ==&macro_url=https%3A%2F%2Fraw.githubusercontent.com%2FdFoiler%2Fnotes%2Fmaster%2Fnir.tex
	\[\begin{tikzcd}[cramped]
		\bullet & \circ & \bullet & \bullet & \circ & \bullet
		\arrow[no head, from=1-1, to=1-2]
		\arrow[no head, from=1-3, to=1-4]
		\arrow[no head, from=1-5, to=1-6]
	\end{tikzcd}\]
	are equivalent. Thus, by iteratively playing the Vogan game on each pair of vertices (first on the long roots and then on the short roots), we get three more Vogan diagrams, as follows.
	\begin{itemize}
		\item We have the following.
		% https://q.uiver.app/#q=WzAsNCxbMCwwLCJcXGJ1bGxldCJdLFsxLDAsIlxcY2lyYyJdLFsyLDAsIlxcY2lyYyJdLFszLDAsIlxcY2lyYyJdLFszLDIsIiIsMCx7InN0eWxlIjp7ImhlYWQiOnsibmFtZSI6Im5vbmUifX19XSxbMiwxLCIiLDAseyJsZXZlbCI6Mn1dLFsxLDAsIiIsMCx7InN0eWxlIjp7ImhlYWQiOnsibmFtZSI6Im5vbmUifX19XV0=&macro_url=https%3A%2F%2Fraw.githubusercontent.com%2FdFoiler%2Fnotes%2Fmaster%2Fnir.tex
		\[\begin{tikzcd}[cramped]
			\bullet & \circ & \circ & \circ
			\arrow[no head, from=1-2, to=1-1]
			\arrow[Rightarrow, from=1-3, to=1-2]
			\arrow[no head, from=1-4, to=1-3]
		\end{tikzcd}\]
		We may call this real form $F_4^1$. In this case, one can see by inspection that $\theta$ acts by $-1$ on $\mf g_\alpha$ if and only if $\alpha\notin\ZZ^4$. It follows that $\mf k=\mf g^\theta$ has exactly the roots from $B_4=\mf{so}(9)$, so $\mf k=\mf{so}(9)$ by comparing ranks. On the remaining roots, one finds that $\mf p$ is given by the spin representation by computing its weight decomposition. Note that this is not the split form because $\op{rank}\mf k\ne\left|\Phi_+\right|=24$, and it is not the compact form because $\theta$ is nontrivial.
		\item We have the following.
		% https://q.uiver.app/#q=WzAsNCxbMCwwLCJcXGJ1bGxldCJdLFsxLDAsIlxcY2lyYyJdLFsyLDAsIlxcY2lyYyJdLFszLDAsIlxcY2lyYyJdLFszLDIsIiIsMCx7InN0eWxlIjp7ImhlYWQiOnsibmFtZSI6Im5vbmUifX19XSxbMiwxLCIiLDAseyJsZXZlbCI6Mn1dLFsxLDAsIiIsMCx7InN0eWxlIjp7ImhlYWQiOnsibmFtZSI6Im5vbmUifX19XV0=&macro_url=https%3A%2F%2Fraw.githubusercontent.com%2FdFoiler%2Fnotes%2Fmaster%2Fnir.tex
		\[\begin{tikzcd}[cramped]
			\bullet & \circ & \circ & \bullet
			\arrow[no head, from=1-2, to=1-1]
			\arrow[Rightarrow, from=1-3, to=1-2]
			\arrow[no head, from=1-4, to=1-3]
		\end{tikzcd}\]
		\item We have the following.
		% https://q.uiver.app/#q=WzAsNCxbMCwwLCJcXGJ1bGxldCJdLFsxLDAsIlxcY2lyYyJdLFsyLDAsIlxcY2lyYyJdLFszLDAsIlxcY2lyYyJdLFszLDIsIiIsMCx7InN0eWxlIjp7ImhlYWQiOnsibmFtZSI6Im5vbmUifX19XSxbMiwxLCIiLDAseyJsZXZlbCI6Mn1dLFsxLDAsIiIsMCx7InN0eWxlIjp7ImhlYWQiOnsibmFtZSI6Im5vbmUifX19XV0=&macro_url=https%3A%2F%2Fraw.githubusercontent.com%2FdFoiler%2Fnotes%2Fmaster%2Fnir.tex
		\[\begin{tikzcd}[cramped]
			\circ & \circ & \circ & \bullet
			\arrow[no head, from=1-2, to=1-1]
			\arrow[Rightarrow, from=1-3, to=1-2]
			\arrow[no head, from=1-4, to=1-3]
		\end{tikzcd}\]
		In fact, this is equivalent to the previous diagram: the previous diagram is equivalent to
		% https://q.uiver.app/#q=WzAsNCxbMCwwLCJcXGJ1bGxldCJdLFsxLDAsIlxcY2lyYyJdLFsyLDAsIlxcY2lyYyJdLFszLDAsIlxcY2lyYyJdLFszLDIsIiIsMCx7InN0eWxlIjp7ImhlYWQiOnsibmFtZSI6Im5vbmUifX19XSxbMiwxLCIiLDAseyJsZXZlbCI6Mn1dLFsxLDAsIiIsMCx7InN0eWxlIjp7ImhlYWQiOnsibmFtZSI6Im5vbmUifX19XV0=&macro_url=https%3A%2F%2Fraw.githubusercontent.com%2FdFoiler%2Fnotes%2Fmaster%2Fnir.tex
		\[\begin{tikzcd}[cramped]
			\circ & \circ & \bullet & \bullet
			\arrow[no head, from=1-2, to=1-1]
			\arrow[Rightarrow, from=1-3, to=1-2]
			\arrow[no head, from=1-4, to=1-3]
		\end{tikzcd}\]
		which is equivalent to
		% https://q.uiver.app/#q=WzAsNCxbMCwwLCJcXGJ1bGxldCJdLFsxLDAsIlxcY2lyYyJdLFsyLDAsIlxcY2lyYyJdLFszLDAsIlxcY2lyYyJdLFszLDIsIiIsMCx7InN0eWxlIjp7ImhlYWQiOnsibmFtZSI6Im5vbmUifX19XSxbMiwxLCIiLDAseyJsZXZlbCI6Mn1dLFsxLDAsIiIsMCx7InN0eWxlIjp7ImhlYWQiOnsibmFtZSI6Im5vbmUifX19XV0=&macro_url=https%3A%2F%2Fraw.githubusercontent.com%2FdFoiler%2Fnotes%2Fmaster%2Fnir.tex
		\[\begin{tikzcd}[cramped]
			\circ & \bullet & \bullet & \circ
			\arrow[no head, from=1-2, to=1-1]
			\arrow[Rightarrow, from=1-3, to=1-2]
			\arrow[no head, from=1-4, to=1-3]
		\end{tikzcd}\]
		which is equivalent to
		\[\begin{tikzcd}[cramped]
			\bullet & \bullet & \bullet & \circ
			\arrow[no head, from=1-2, to=1-1]
			\arrow[Rightarrow, from=1-3, to=1-2]
			\arrow[no head, from=1-4, to=1-3]
		\end{tikzcd}\]
		which is equivalent to
		\[\begin{tikzcd}[cramped]
			\circ & \circ & \bullet & \bullet
			\arrow[no head, from=1-2, to=1-1]
			\arrow[Rightarrow, from=1-3, to=1-2]
			\arrow[no head, from=1-4, to=1-3]
		\end{tikzcd}\]
		which is equivalent to
		\[\begin{tikzcd}[cramped]
			\circ & \circ & \circ & \bullet
			\arrow[no head, from=1-2, to=1-1]
			\arrow[Rightarrow, from=1-3, to=1-2]
			\arrow[no head, from=1-4, to=1-3]
		\end{tikzcd}\]
		as needed.

		Now, by exhaustion, we conclude that we must have the split form. Here, we find that $\theta$ fixes a root system of type $C_3$, so $\mf k$ contains $\mf{sp}(6)$, which has dimension $21$. However, $\mf k$ should increase in rank by $1$ while increasing in dimension $3$, so we are required to have $\mf k=\mf{sp}(6)\oplus\mf{sl}(2)$. Indeed, all semisimple algebras strictly containing $\mf{sp}(6)$ have dimension larger than $24$, so $\mf k$ must have another summand, and the smallest possible summand to increase the rank is either $\mf{sl}(2)$ or $\mf{gl}(1)$, but adding $\mf{gl}(1)$s cannot help us.
		\qedhere
	\end{itemize}
\end{proof}

\subsection{Real Forms of Type \texorpdfstring{$E$}{ E}}
Having done everything easy, it remains to do things which are hard. We are now in type $E$. To start, note that these types have a single inner class, except the Dynkin diagram $E_6$ admits a nontrivial automorphism. Let's deal with this first.
\begin{exe}
	We compute the real forms in type $E_6$ which are in the non-compact inner class.
\end{exe}
\begin{proof}
	In type $E_6$, it is possible to have vertices which are connected in the Vogan diagram: indeed, the split inner class comes from the flip of the Dynkin diagram. In this case, playing the Vogan diagram on the remaining two fixed vertices, we get two real forms, as follows.
	\begin{itemize}
		\item We have the following.
		% https://q.uiver.app/#q=WzAsNixbMCwwLCJcXGNpcmMiXSxbMSwwLCJcXGNpcmMiXSxbMiwwLCJcXGNpcmMiXSxbMywwLCJcXGNpcmMiXSxbNCwwLCJcXGNpcmMiXSxbMiwxLCJcXGNpcmMiXSxbMiw1LCIiLDAseyJzdHlsZSI6eyJoZWFkIjp7Im5hbWUiOiJub25lIn19fV0sWzAsMSwiIiwwLHsic3R5bGUiOnsiaGVhZCI6eyJuYW1lIjoibm9uZSJ9fX1dLFsxLDIsIiIsMCx7InN0eWxlIjp7ImhlYWQiOnsibmFtZSI6Im5vbmUifX19XSxbMiwzLCIiLDAseyJzdHlsZSI6eyJoZWFkIjp7Im5hbWUiOiJub25lIn19fV0sWzMsNCwiIiwwLHsic3R5bGUiOnsiaGVhZCI6eyJuYW1lIjoibm9uZSJ9fX1dLFsxLDMsIiIsMSx7ImN1cnZlIjotMSwic3R5bGUiOnsiaGVhZCI6eyJuYW1lIjoibm9uZSJ9fX1dLFswLDQsIiIsMix7ImN1cnZlIjotMiwic3R5bGUiOnsiaGVhZCI6eyJuYW1lIjoibm9uZSJ9fX1dXQ==&macro_url=https%3A%2F%2Fraw.githubusercontent.com%2FdFoiler%2Fnotes%2Fmaster%2Fnir.tex
		\[\begin{tikzcd}[cramped]
			\circ & \circ & \circ & \circ & \circ \\
			&& \circ
			\arrow[no head, from=1-1, to=1-2]
			\arrow[curve={height=-12pt}, no head, from=1-1, to=1-5]
			\arrow[no head, from=1-2, to=1-3]
			\arrow[curve={height=-6pt}, no head, from=1-2, to=1-4]
			\arrow[no head, from=1-3, to=1-4]
			\arrow[no head, from=1-3, to=2-3]
			\arrow[no head, from=1-4, to=1-5]
		\end{tikzcd}\]
		Here, we see that $\mf k$ has rank $4$, and by looking at the projection of these roots to $\mf k$, we see that the two pairs of associated roots become short, so we $\mf k$ has type $F_4$. This is not the split form for dimension reasons: here, $\dim E_6=78$, and the split form has $\dim\mf k$ equal to the number of positive roots.
		\item The Vogan diagram
		% https://q.uiver.app/#q=WzAsNixbMCwwLCJcXGNpcmMiXSxbMSwwLCJcXGNpcmMiXSxbMiwwLCJcXGNpcmMiXSxbMywwLCJcXGNpcmMiXSxbNCwwLCJcXGNpcmMiXSxbMiwxLCJcXGJ1bGxldCJdLFsyLDUsIiIsMCx7InN0eWxlIjp7ImhlYWQiOnsibmFtZSI6Im5vbmUifX19XSxbMCwxLCIiLDAseyJzdHlsZSI6eyJoZWFkIjp7Im5hbWUiOiJub25lIn19fV0sWzEsMiwiIiwwLHsic3R5bGUiOnsiaGVhZCI6eyJuYW1lIjoibm9uZSJ9fX1dLFsyLDMsIiIsMCx7InN0eWxlIjp7ImhlYWQiOnsibmFtZSI6Im5vbmUifX19XSxbMyw0LCIiLDAseyJzdHlsZSI6eyJoZWFkIjp7Im5hbWUiOiJub25lIn19fV0sWzEsMywiIiwxLHsiY3VydmUiOi0xLCJzdHlsZSI6eyJoZWFkIjp7Im5hbWUiOiJub25lIn19fV0sWzAsNCwiIiwyLHsiY3VydmUiOi0yLCJzdHlsZSI6eyJoZWFkIjp7Im5hbWUiOiJub25lIn19fV1d&macro_url=https%3A%2F%2Fraw.githubusercontent.com%2FdFoiler%2Fnotes%2Fmaster%2Fnir.tex
		\[\begin{tikzcd}[cramped]
			\circ & \circ & \circ & \circ & \circ \\
			&& \bullet
			\arrow[no head, from=1-1, to=1-2]
			\arrow[curve={height=-12pt}, no head, from=1-1, to=1-5]
			\arrow[no head, from=1-2, to=1-3]
			\arrow[curve={height=-6pt}, no head, from=1-2, to=1-4]
			\arrow[no head, from=1-3, to=1-4]
			\arrow[no head, from=1-3, to=2-3]
			\arrow[no head, from=1-4, to=1-5]
		\end{tikzcd}\]
		must correspond to the split form for dimension reasons. The same projection argument shows that $\mf k$ contains $\mf{sp}(6)$. One can check that it is $\mf{sp}(8)$, which we will do on the homework.
		\qedhere
	\end{itemize}
\end{proof}
We now turn to the compact inner class. Let's start by making some general comments. The Dynkin diagram has a short leg, a middle leg, and a long leg. For example, $E_7$ looks like the following.
% https://q.uiver.app/#q=WzAsNyxbMCwwLCJcXGNpcmMiXSxbMSwwLCJcXGNpcmMiXSxbMiwwLCJcXGNpcmMiXSxbMiwxLCJcXGNpcmMiXSxbMywwLCJcXGNpcmMiXSxbNCwwLCJcXGNpcmMiXSxbNSwwLCJcXGNpcmMiXSxbMCwxLCIiLDAseyJzdHlsZSI6eyJoZWFkIjp7Im5hbWUiOiJub25lIn19fV0sWzEsMiwiIiwwLHsic3R5bGUiOnsiaGVhZCI6eyJuYW1lIjoibm9uZSJ9fX1dLFsyLDQsIiIsMCx7InN0eWxlIjp7ImhlYWQiOnsibmFtZSI6Im5vbmUifX19XSxbNCw1LCIiLDAseyJzdHlsZSI6eyJoZWFkIjp7Im5hbWUiOiJub25lIn19fV0sWzUsNiwiIiwwLHsic3R5bGUiOnsiaGVhZCI6eyJuYW1lIjoibm9uZSJ9fX1dLFszLDIsIiIsMCx7InN0eWxlIjp7ImhlYWQiOnsibmFtZSI6Im5vbmUifX19XV0=&macro_url=https%3A%2F%2Fraw.githubusercontent.com%2FdFoiler%2Fnotes%2Fmaster%2Fnir.tex
\[\begin{tikzcd}[cramped]
	\circ & \circ & \circ & \circ & \circ & \circ \\
	&& \circ
	\arrow[no head, from=1-1, to=1-2]
	\arrow[no head, from=1-2, to=1-3]
	\arrow[no head, from=1-3, to=1-4]
	\arrow[no head, from=1-4, to=1-5]
	\arrow[no head, from=1-5, to=1-6]
	\arrow[no head, from=2-3, to=1-3]
\end{tikzcd}\]
Once we exclude the compact form, we may assume that there is at least one black vertex. By playing the Vogan game, we may move the black vertex to the end of the short leg, and then we can always use this short leg in order to force the intersection point to be a black vertex. We now do casework on our type.

We now use this to force the long leg of the Dynkin diagram into a particular shape, possibly at the cost of the black vertex at the short leg.
\begin{itemize}
	\item In type $E_6$, the short and long leg has two vertices, so there are two possibilities on each leg: $\circ-\circ$ or $\circ-\bullet$. By possibly undoing our work with the nodal vertex, we may force the long leg to be $\bullet-\bullet$.
	\item In type $E_7$, we also have to play the Vogan diagram on three vertices. We find that the long leg has $\circ-\circ-\circ$ or $\bullet-\circ-\circ=\bullet-\bullet-\circ=\circ-\bullet-\bullet$ or $\bullet-\circ-\bullet=\bullet-\bullet-\bullet=\circ-\bullet-\circ$. However, using the nodal vertex allows us to see that the rightmost diagrams of the last two classes are equivalent, and the leftmost diagrams of the last two classes are equivalent.
	\item In type $E_8$, the long leg has four vertices. One again plays the Vogan game on $A_4$. We have the following classes.
	\begin{itemize}
		\item We have $\circ-\circ-\circ-\circ$.
		\item We have $\bullet-\circ-\circ-\circ=\bullet-\bullet-\circ-\circ=\circ-\bullet-\bullet-\circ=\circ-\circ-\bullet-\bullet=\circ-\circ-\circ-\bullet$. Note that the nodal vertex allows us to push this to the previous class.
		\item We have $\bullet-\circ-\circ-\bullet=\bullet-\circ-\bullet-\bullet=\bullet-\bullet-\bullet-\circ=\circ-\bullet-\circ-\circ$ or $\bullet-\circ-\bullet-\circ=\bullet-\bullet-\bullet-\bullet=\bullet-\bullet-\circ-\bullet=\circ-\bullet-\bullet-\bullet$. Again, the nodal vertex allows us to push this to $\circ-\circ-\circ-\bullet$ and thus to the previous class.
	\end{itemize}
\end{itemize}
Now, by redoing our nodal vertex, we find that we may assume that the short leg can be assumed to be black, the node can be black, and the long leg can look like $?-\bullet-\cdots-\bullet$. In types $E_6$ or $E_8$, one finds that we can even assume that $?$ is black by some combinatorics: in $E_6$, we are simply saying that $\cdots-\circ-\bullet=\cdots-\bullet-\bullet$. We won't write it our for $E_8$.

Similarly, the middle leg is a Vogan game on $A_2$, so it has either $\circ-\circ$ or $\bullet-\bullet=\bullet-\circ=\circ-\bullet$. Thus, $E_6$ has at most two Vogan diagrams, $E_7$ has at most four Vogan diagrams, and $E_8$ has at most two Vogan diagrams. We now work through our types separately.
\begin{exe}
	We compute the real forms in type $E_6$ in the compact inner class.
\end{exe}
\begin{proof}
	We know that there are at most three Vogan diagrams in the compact inner class, two of which are non-compact. Here are our diagrams.
	\begin{itemize}
		\item We have the following.
		% https://q.uiver.app/#q=WzAsNixbMCwwLCJcXGNpcmMiXSxbMSwwLCJcXGNpcmMiXSxbMiwwLCJcXGNpcmMiXSxbMiwxLCJcXGNpcmMiXSxbMywwLCJcXGNpcmMiXSxbNCwwLCJcXGJ1bGxldCJdLFswLDEsIiIsMCx7InN0eWxlIjp7ImhlYWQiOnsibmFtZSI6Im5vbmUifX19XSxbMSwyLCIiLDAseyJzdHlsZSI6eyJoZWFkIjp7Im5hbWUiOiJub25lIn19fV0sWzIsNCwiIiwwLHsic3R5bGUiOnsiaGVhZCI6eyJuYW1lIjoibm9uZSJ9fX1dLFs0LDUsIiIsMCx7InN0eWxlIjp7ImhlYWQiOnsibmFtZSI6Im5vbmUifX19XSxbMiwzLCIiLDAseyJzdHlsZSI6eyJoZWFkIjp7Im5hbWUiOiJub25lIn19fV1d&macro_url=https%3A%2F%2Fraw.githubusercontent.com%2FdFoiler%2Fnotes%2Fmaster%2Fnir.tex
		\[\begin{tikzcd}[cramped]
			\circ & \circ & \circ & \circ & \bullet \\
			&& \circ
			\arrow[no head, from=1-1, to=1-2]
			\arrow[no head, from=1-2, to=1-3]
			\arrow[no head, from=1-3, to=1-4]
			\arrow[no head, from=1-3, to=2-3]
			\arrow[no head, from=1-4, to=1-5]
		\end{tikzcd}\]
		Here, we label the black vertex to be the miniscule weight $\alpha_1=\omega_1$. It follows that every positive root contains $\alpha_1$ with a coefficient of either $0$ or $1$ (because it is a fundamental weight), and $\theta$ fixes the root if and only if this coefficient is $+1$. We conclude that $\mf k$ contains a copy of the rest of the Dynkin diagram, which gives $\mf{so}(5)\subseteq\mf k$. However, the rank of $\mf k$ should be five, and a dimension argument (with the signature of the Killing form) gives $\mf k=\mf{so}(5)\oplus\mf{so}(2)$.
		\item We have the following.
		% https://q.uiver.app/#q=WzAsNixbMCwwLCJcXGNpcmMiXSxbMSwwLCJcXGNpcmMiXSxbMiwwLCJcXGNpcmMiXSxbMiwxLCJcXGJ1bGxldCJdLFszLDAsIlxcY2lyYyJdLFs0LDAsIlxcY2lyYyJdLFswLDEsIiIsMCx7InN0eWxlIjp7ImhlYWQiOnsibmFtZSI6Im5vbmUifX19XSxbMSwyLCIiLDAseyJzdHlsZSI6eyJoZWFkIjp7Im5hbWUiOiJub25lIn19fV0sWzIsNCwiIiwwLHsic3R5bGUiOnsiaGVhZCI6eyJuYW1lIjoibm9uZSJ9fX1dLFs0LDUsIiIsMCx7InN0eWxlIjp7ImhlYWQiOnsibmFtZSI6Im5vbmUifX19XSxbMiwzLCIiLDAseyJzdHlsZSI6eyJoZWFkIjp7Im5hbWUiOiJub25lIn19fV1d&macro_url=https%3A%2F%2Fraw.githubusercontent.com%2FdFoiler%2Fnotes%2Fmaster%2Fnir.tex
		\[\begin{tikzcd}[cramped]
			\circ & \circ & \circ & \circ & \circ \\
			&& \bullet
			\arrow[no head, from=1-1, to=1-2]
			\arrow[no head, from=1-2, to=1-3]
			\arrow[no head, from=1-3, to=1-4]
			\arrow[no head, from=1-3, to=2-3]
			\arrow[no head, from=1-4, to=1-5]
		\end{tikzcd}\]
		The same argument shows that $\mf k$ contains $\mf{sl}(6)$. The black vertex corresponds to the fundamental weight of the adjoint representation, and one can check that this means that $\theta$ will act by $+1$ only on an extra copy of $\mf{sl}(2)$ given by the root spaces of the maximal root. (Notably, $\mf k$ exactly contains the torus and the roots where $\theta$ acts by $+1$, which can be computed.)
		\qedhere
	\end{itemize}
\begin{exe}
	We compute the real forms in type $E_7$ in the compact inner class.
\end{exe}
\begin{proof}
	There are at most four Vogan diagrams which have a black vertex, though we will eventually bring that number down to three. Note $\dim E_7=133$, so the number of positive roots is $63$.
	\begin{itemize}
		\item We have the following.
		% https://q.uiver.app/#q=WzAsNyxbMCwwLCJcXGNpcmMiXSxbMSwwLCJcXGNpcmMiXSxbMiwwLCJcXGNpcmMiXSxbMiwxLCJcXGNpcmMiXSxbMywwLCJcXGNpcmMiXSxbNCwwLCJcXGNpcmMiXSxbNSwwLCJcXGJ1bGxldCJdLFswLDEsIiIsMCx7InN0eWxlIjp7ImhlYWQiOnsibmFtZSI6Im5vbmUifX19XSxbMSwyLCIiLDAseyJzdHlsZSI6eyJoZWFkIjp7Im5hbWUiOiJub25lIn19fV0sWzIsNCwiIiwwLHsic3R5bGUiOnsiaGVhZCI6eyJuYW1lIjoibm9uZSJ9fX1dLFs0LDUsIiIsMCx7InN0eWxlIjp7ImhlYWQiOnsibmFtZSI6Im5vbmUifX19XSxbMiwzLCIiLDAseyJzdHlsZSI6eyJoZWFkIjp7Im5hbWUiOiJub25lIn19fV0sWzYsNSwiIiwyLHsic3R5bGUiOnsiaGVhZCI6eyJuYW1lIjoibm9uZSJ9fX1dXQ==&macro_url=https%3A%2F%2Fraw.githubusercontent.com%2FdFoiler%2Fnotes%2Fmaster%2Fnir.tex
		\[\begin{tikzcd}[cramped]
			\circ & \circ & \circ & \circ & \circ & \bullet \\
			&& \circ
			\arrow[no head, from=1-1, to=1-2]
			\arrow[no head, from=1-2, to=1-3]
			\arrow[no head, from=1-3, to=1-4]
			\arrow[no head, from=1-3, to=2-3]
			\arrow[no head, from=1-4, to=1-5]
			\arrow[no head, from=1-6, to=1-5]
		\end{tikzcd}\]
		Here, $\mf k$ has a copy of $E_6$, which has dimension $78$, which is bigger than $63$, so $\mf k$ is not split. Because the colored root is miniscule, we find that $\mf k=E_6\oplus\mf{so}(2)$ again.
		\item We have the following.
		\[\begin{tikzcd}[cramped]
			\bullet & \circ & \circ & \circ & \circ & \circ \\
			&& \circ
			\arrow[no head, from=1-1, to=1-2]
			\arrow[no head, from=1-2, to=1-3]
			\arrow[no head, from=1-3, to=1-4]
			\arrow[no head, from=1-3, to=2-3]
			\arrow[no head, from=1-4, to=1-5]
			\arrow[no head, from=1-6, to=1-5]
		\end{tikzcd}\]
		Here, $\mf k$ contains $\mf{so}(12)$ automatically, so dimension reasons make this form not split, and we can see that $\mf k=\mf{so}(12)\oplus\mf{sl}(2)$ because the black vertex is maximal, corresponding to the adjoint representation.
		\item We have the following.
		\[\begin{tikzcd}[cramped]
			\circ & \circ & \circ & \circ & \circ & \circ \\
			&& \bullet
			\arrow[no head, from=1-1, to=1-2]
			\arrow[no head, from=1-2, to=1-3]
			\arrow[no head, from=1-3, to=1-4]
			\arrow[no head, from=1-3, to=2-3]
			\arrow[no head, from=1-4, to=1-5]
			\arrow[no head, from=1-6, to=1-5]
		\end{tikzcd}\]
		Here, $\mf k$ contains $\mf{sl}(7)$. This will turn out to be the split form by exhaustion, so it has dimension 
		\item It remains to show that our last two Vogan diagrams were actually the same. Well,
		% https://q.uiver.app/#q=WzAsNyxbMCwwLCJcXGJ1bGxldCJdLFsxLDAsIlxcY2lyYyJdLFsyLDAsIlxcYnVsbGV0Il0sWzMsMCwiXFxjaXJjIl0sWzQsMCwiXFxidWxsZXQiXSxbNSwwLCJcXGJ1bGxldCJdLFsyLDEsIlxcYnVsbGV0Il0sWzAsMSwiIiwwLHsic3R5bGUiOnsiaGVhZCI6eyJuYW1lIjoibm9uZSJ9fX1dLFsxLDIsIiIsMCx7InN0eWxlIjp7ImhlYWQiOnsibmFtZSI6Im5vbmUifX19XSxbMiwzLCIiLDAseyJzdHlsZSI6eyJoZWFkIjp7Im5hbWUiOiJub25lIn19fV0sWzMsNCwiIiwwLHsic3R5bGUiOnsiaGVhZCI6eyJuYW1lIjoibm9uZSJ9fX1dLFs0LDUsIiIsMCx7InN0eWxlIjp7ImhlYWQiOnsibmFtZSI6Im5vbmUifX19XSxbMiw2LCIiLDAseyJzdHlsZSI6eyJoZWFkIjp7Im5hbWUiOiJub25lIn19fV1d&macro_url=https%3A%2F%2Fraw.githubusercontent.com%2FdFoiler%2Fnotes%2Fmaster%2Fnir.tex
		\[\begin{tikzcd}[cramped]
			\bullet & \circ & \bullet & \circ & \bullet & \bullet \\
			&& \bullet
			\arrow[no head, from=1-1, to=1-2]
			\arrow[no head, from=1-2, to=1-3]
			\arrow[no head, from=1-3, to=1-4]
			\arrow[no head, from=1-3, to=2-3]
			\arrow[no head, from=1-4, to=1-5]
			\arrow[no head, from=1-5, to=1-6]
		\end{tikzcd}\]
		is
		\[\begin{tikzcd}[cramped]
			\bullet & \bullet & \bullet & \bullet & \bullet & \bullet \\
			&& \bullet
			\arrow[no head, from=1-1, to=1-2]
			\arrow[no head, from=1-2, to=1-3]
			\arrow[no head, from=1-3, to=1-4]
			\arrow[no head, from=1-3, to=2-3]
			\arrow[no head, from=1-4, to=1-5]
			\arrow[no head, from=1-5, to=1-6]
		\end{tikzcd}\]
		which is
		\[\begin{tikzcd}[cramped]
			\circ & \bullet & \circ & \bullet & \bullet & \bullet \\
			&& \bullet
			\arrow[no head, from=1-1, to=1-2]
			\arrow[no head, from=1-2, to=1-3]
			\arrow[no head, from=1-3, to=1-4]
			\arrow[no head, from=1-3, to=2-3]
			\arrow[no head, from=1-4, to=1-5]
			\arrow[no head, from=1-5, to=1-6]
		\end{tikzcd}\]
		which is
		\[\begin{tikzcd}[cramped]
			\circ & \bullet & \bullet & \bullet & \circ & \bullet \\
			&& \bullet
			\arrow[no head, from=1-1, to=1-2]
			\arrow[no head, from=1-2, to=1-3]
			\arrow[no head, from=1-3, to=1-4]
			\arrow[no head, from=1-3, to=2-3]
			\arrow[no head, from=1-4, to=1-5]
			\arrow[no head, from=1-5, to=1-6]
		\end{tikzcd}\]
		which is
		\[\begin{tikzcd}[cramped]
			\circ & \circ & \bullet & \circ & \circ & \circ \\
			&& \bullet
			\arrow[no head, from=1-1, to=1-2]
			\arrow[no head, from=1-2, to=1-3]
			\arrow[no head, from=1-3, to=1-4]
			\arrow[no head, from=1-3, to=2-3]
			\arrow[no head, from=1-4, to=1-5]
			\arrow[no head, from=1-5, to=1-6]
		\end{tikzcd}\]
		which is
		\[\begin{tikzcd}[cramped]
			\circ & \circ & \bullet & \circ & \circ & \bullet \\
			&& \bullet
			\arrow[no head, from=1-1, to=1-2]
			\arrow[no head, from=1-2, to=1-3]
			\arrow[no head, from=1-3, to=1-4]
			\arrow[no head, from=1-3, to=2-3]
			\arrow[no head, from=1-4, to=1-5]
			\arrow[no head, from=1-5, to=1-6]
		\end{tikzcd}\]
		which is among one of our restricted diagrams.
		\qedhere
	\end{itemize}
\end{proof}
\begin{exe}
	We compute the real forms in type $E_8$ in the compact inner class.
\end{exe}
\begin{proof}
	Note $E_8$ has $120$ positive roots.
	\begin{itemize}
		\item Consider the following Vogan diagram.
		% https://q.uiver.app/#q=WzAsOCxbMCwwLCJcXGNpcmMiXSxbMSwwLCJcXGNpcmMiXSxbMiwwLCJcXGNpcmMiXSxbMiwxLCJcXGNpcmMiXSxbMywwLCJcXGNpcmMiXSxbNCwwLCJcXGNpcmMiXSxbNSwwLCJcXGNpcmMiXSxbNiwwLCJcXGJ1bGxldCJdLFswLDEsIiIsMCx7InN0eWxlIjp7ImhlYWQiOnsibmFtZSI6Im5vbmUifX19XSxbMSwyLCIiLDAseyJzdHlsZSI6eyJoZWFkIjp7Im5hbWUiOiJub25lIn19fV0sWzIsNCwiIiwwLHsic3R5bGUiOnsiaGVhZCI6eyJuYW1lIjoibm9uZSJ9fX1dLFs0LDUsIiIsMCx7InN0eWxlIjp7ImhlYWQiOnsibmFtZSI6Im5vbmUifX19XSxbMiwzLCIiLDAseyJzdHlsZSI6eyJoZWFkIjp7Im5hbWUiOiJub25lIn19fV0sWzYsNSwiIiwyLHsic3R5bGUiOnsiaGVhZCI6eyJuYW1lIjoibm9uZSJ9fX1dLFs3LDYsIiIsMix7InN0eWxlIjp7ImhlYWQiOnsibmFtZSI6Im5vbmUifX19XV0=&macro_url=https%3A%2F%2Fraw.githubusercontent.com%2FdFoiler%2Fnotes%2Fmaster%2Fnir.tex
		\[\begin{tikzcd}[cramped]
			\circ & \circ & \circ & \circ & \circ & \circ & \bullet \\
			&& \circ
			\arrow[no head, from=1-1, to=1-2]
			\arrow[no head, from=1-2, to=1-3]
			\arrow[no head, from=1-3, to=1-4]
			\arrow[no head, from=1-3, to=2-3]
			\arrow[no head, from=1-4, to=1-5]
			\arrow[no head, from=1-6, to=1-5]
			\arrow[no head, from=1-7, to=1-6]
		\end{tikzcd}\]
		Here, the black vertex is the adjoint representation, so $\mf k=E_7\oplus\mf{sl}(2)$ for rank reasons. This is not the split form because $\dim\mf k$ is too large.
		\item Consider the following Vogan diagram.
		% https://q.uiver.app/#q=WzAsOCxbMCwwLCJcXGJ1bGxldCJdLFsxLDAsIlxcY2lyYyJdLFsyLDAsIlxcY2lyYyJdLFsyLDEsIlxcY2lyYyJdLFszLDAsIlxcY2lyYyJdLFs0LDAsIlxcY2lyYyJdLFs1LDAsIlxcY2lyYyJdLFs2LDAsIlxcY2lyYyJdLFswLDEsIiIsMCx7InN0eWxlIjp7ImhlYWQiOnsibmFtZSI6Im5vbmUifX19XSxbMSwyLCIiLDAseyJzdHlsZSI6eyJoZWFkIjp7Im5hbWUiOiJub25lIn19fV0sWzIsNCwiIiwwLHsic3R5bGUiOnsiaGVhZCI6eyJuYW1lIjoibm9uZSJ9fX1dLFs0LDUsIiIsMCx7InN0eWxlIjp7ImhlYWQiOnsibmFtZSI6Im5vbmUifX19XSxbMiwzLCIiLDAseyJzdHlsZSI6eyJoZWFkIjp7Im5hbWUiOiJub25lIn19fV0sWzYsNSwiIiwyLHsic3R5bGUiOnsiaGVhZCI6eyJuYW1lIjoibm9uZSJ9fX1dLFs3LDYsIiIsMix7InN0eWxlIjp7ImhlYWQiOnsibmFtZSI6Im5vbmUifX19XV0=&macro_url=https%3A%2F%2Fraw.githubusercontent.com%2FdFoiler%2Fnotes%2Fmaster%2Fnir.tex
		\[\begin{tikzcd}[cramped]
			\bullet & \circ & \circ & \circ & \circ & \circ & \circ \\
			&& \circ
			\arrow[no head, from=1-1, to=1-2]
			\arrow[no head, from=1-2, to=1-3]
			\arrow[no head, from=1-3, to=1-4]
			\arrow[no head, from=1-3, to=2-3]
			\arrow[no head, from=1-4, to=1-5]
			\arrow[no head, from=1-6, to=1-5]
			\arrow[no head, from=1-7, to=1-6]
		\end{tikzcd}\]
		Here, $\mf k$ contains $\mf{so}(14)$. This should be the split form by exhaustion. For dimension reasons, this needs to be $\mf{so}(16)$. Also, $\mf k$ needs to look like $\mf{so}(14)\oplus R$ where $R$ is some remaining representation of $\mf{so}(14)$ of dimension $29$, which we see must be $\CC^{14}\oplus\CC^{14}\oplus\CC$. (This representation needs to be self-dual.) By inspection, we must have $\mf{so}(16)$. While we're here, we remark that $\mf p$ turns out to be the spin representation.
		\qedhere
	\end{itemize}
\end{proof}

\end{document}